\chapter{Industry \& Competitive Analysis}\label{ch:competitive-analysis}

% Scope: the five forces, the value chain, industry structure and profitability (HHI).
% Cross-link: durable advantage that survives these forces is the subject of
%   \cref{ch:advantage-and-disruption}; price power traced here roots in
%   \cref{ch:pricing} (Lerner index) and \cref{ch:customer-economics} (LTV).

Industry profitability is not random. The structural conditions that surround a firm
--- supplier leverage, buyer leverage, substitutes, new entrants, and internal rivalry ---
set a ceiling on returns before strategy even begins. The value chain then asks where,
within that industry, the firm captures what remains.

\section{Porter's Five Forces}\label{sec:five-forces}
% complexity: 4 -> figure five-forces

Before allocating capital or entering a market, a strategist needs to know which structural
features will cap returns and which will support them. Five Forces answers that question before
a single customer is signed. The mechanism is straightforward: a high-margin market attracts
entrants and powerful buyers until returns compress to the cost of capital. Five Forces names
the five channels through which that compression happens and asks how fierce each is right now
--- not whether the industry is attractive in the abstract, but whether \emph{this firm} can
hold a position in it.

The framework is qualitative; its analytical work is a diagnostic sequence applied to each
force in turn \parencite{porter1980}:
\begin{description}
  \item[Rivalry among existing competitors] Concentration (see \cref{eq:hhi}),
        price transparency, product differentiation, exit costs, and capacity increments.
        High fixed costs and commodity products make rivalry intense.
  \item[Threat of new entrants] Minimum efficient scale, proprietary technology,
        network effects, brand loyalty, switching costs, and regulatory barriers.
        Low capital requirements and mature technology lower this barrier.
  \item[Threat of substitutes] Products from adjacent categories that fulfil the same
        job-to-be-done; assessed by cross-price elasticity (\cref{eq:elasticity}).
        Substitutes set a soft price ceiling.
  \item[Bargaining power of buyers] Buyer concentration, switching costs, backward-
        integration threat, share of buyer's spend, and price sensitivity.
  \item[Bargaining power of suppliers] Supplier concentration, uniqueness of input,
        forward-integration threat, and switching cost of changing supplier.
\end{description}

\begin{figure}[htbp]
  \centering
  \includegraphics[width=0.86\linewidth]{five-forces}
  \caption{The Five Forces. Rivalry sits at the centre; four forces exert inward pressure on
    industry profitability. Strength of each arrow represents the aggregate bargaining or
    threat intensity.}
  \label{fig:five-forces}
\end{figure}

\begin{admonition}{Warning}
Five Forces is a static, industry-level snapshot. It misses three things that matter
increasingly in technology markets. First, \emph{platform complements}: a marketplace's
profitability depends on complementor investment, which is neither supplier nor buyer in the
traditional sense. Second, \emph{dynamic competition}: the model has no clock; a force that
looks weak today (e.g.\ entry) may be irrelevant in a rapidly restructuring market.
Third, the resource-based view (\cref{ch:advantage-and-disruption}) argues that
firm-level capability heterogeneity --- not industry structure alone --- explains return
differences inside an industry. Use Five Forces to pick the arena; use VRIN to ask whether
you can win in it.
\end{admonition}

\begin{example}
Apply to the running B2B SaaS market. \emph{Decision: should we invest in differentiation
or compete on price?}

\begin{itemize}
  \item \textbf{Buyer power --- high.} Buyers are SMBs with low switching costs (data exports
        are standard), multiple alternatives, and low per-seat spend relative to their budgets.
        This compresses price power toward the Lerner index floor (\cref{eq:lerner}).
  \item \textbf{Entry threat --- high.} No-code infrastructure and open APIs mean a credible
        point solution can be assembled for well under \money{500000}. Network effects are
        weak at launch.
  \item \textbf{Rivalry --- intense.} The market is fragmented (HHI $\approx$\num{2150},
        see \cref{eq:hhi}); feature parity is common; price comparison sites make switching
        visible. At \money{50}/month the product is fully in the commodity danger zone.
  \item \textbf{Substitutes --- moderate.} Spreadsheets and lighter SaaS tools serve the
        same job; threat is real but manageable if integrations create stickiness.
  \item \textbf{Supplier power --- low.} Cloud infrastructure (AWS, GCP) is competitive;
        no proprietary input is captive to one vendor.
\end{itemize}

High buyer power plus high entry threat in a fragmented market forces a differentiation
response --- product uniqueness that raises switching costs --- rather than a cost-leadership
race the firm cannot win at its scale.
\end{example}

\begin{admonition}{Warning}
Treating the five forces as equally weighted and averaging them into a single
``attractiveness score.'' A market can be \emph{structurally unattractive} on four
dimensions yet highly profitable for a firm that has neutralised the fifth (e.g.\ a patent
wall erasing entry threat). Score each force separately; then ask which one your strategy
addresses.
\end{admonition}

\textcite{dranove7theconomics} extend Five Forces with game-theoretic analysis of rivalry:
when competitors interact repeatedly, tit-for-tat strategies can sustain supra-competitive
pricing without collusion. The model also sits inside a broader ``structure--conduct--
performance'' (SCP) paradigm whose empirical validity is debated (see
\cref{ch:advantage-and-disruption} for the resource-based critique).

The Five Forces diagnosis sets the context for every pricing decision in
\cref{ch:pricing} and every segmentation choice in
\cref{ch:segmentation-targeting-positioning}: enter only where at least one force is
genuinely manageable. \textcite{porter1980} and \textcite{dranove7theconomics} develop the
empirical industry-economics underpinning.

\section{The Value Chain}\label{sec:value-chain}
% complexity: 3

Which activities inside the firm generate the cost or differentiation that matters to the
customer, and which are overhead that should be minimised or outsourced? A firm is not just
a product; it is a sequence of activities. Cost and differentiation both originate in specific
links of that chain. Knowing which link matters is the precondition for allocating investment
correctly --- and for understanding where a competitor with a different chain structure will
undercut or outperform you.

\textcite{porter1985} organises firm activities into two tiers:
\begin{description}
  \item[Primary activities] Those that directly touch the product or service:
    inbound logistics $\to$ operations $\to$ outbound logistics $\to$
    marketing \& sales $\to$ service.
  \item[Support activities] Those that enable primaries: firm infrastructure,
    human-resource management, technology development, and procurement.
\end{description}
For a B2B SaaS firm the primary chain compresses to four links:
product development $\to$ sales \& marketing $\to$ customer success $\to$
infrastructure/operations. Support activities (finance, legal, HR) apply unchanged.
``Margin'' is the aggregate surplus left after the combined cost of all links.

The value chain is \emph{not} a supply chain: supply chains trace physical inputs and
outputs across firms; the value chain maps value-adding activities \emph{within} a single
firm. Conflating them leads to outsourcing decisions based on cost alone, ignoring the
differentiation and tacit-knowledge embedded in the activity. The model also assumes
activities are separable, which overstates in highly integrated digital products.

\begin{example}
\emph{Decision: where is the SaaS firm spending and where is it building advantage?}

Observed operating expense mix (approximation, consistent with gross margin
$\approx$\qty{70}{\percent}):
\begin{itemize}
  \item Sales \& marketing: $\approx$\qty{50}{\percent} of revenue
        (the dominant link --- consistent with a \money{120000}/month ad budget
        plus a sales team).
  \item R\&D / product: $\approx$\qty{30}{\percent} of revenue.
  \item Customer success / infrastructure: $\approx$\qty{15}{\percent} of revenue.
  \item G\&A: $\approx$\qty{5}{\percent}.
\end{itemize}
S\&M dominates because the firm is pre-moat: it must buy customers it cannot yet retain
organically. A post-moat firm (\cref{ch:growth}) shifts spend toward product and CS as
switching costs reduce the marginal cost of retention. The value-chain lens makes this
trajectory visible as a shift in which link carries the investment.
\end{example}

\begin{admonition}{Warning}
Equating the value chain with an activity cost breakdown and then ``optimising'' each link
independently. The chain's value comes from \emph{linkages}: a superior service activity
raises the value of the product activity by reducing churn (see \cref{eq:retention} in
\cref{ch:customer-economics}). Cutting CS to reduce cost-per-link destroys the linkage
value --- and lifts churn.
\end{admonition}

The value chain predates the platform-business model, where the primary activities are
reversed: the platform produces a matching infrastructure, and value is created in the
interactions between third-party participants rather than in the firm's own operations.
The implications for cost structure and competitive analysis are explored in
\cref{ch:business-models}.

The value chain's cost concentration in S\&M reflects the competitive structure Five Forces
diagnosed above: when entry threat and buyer power are both high, demand must be acquired
expensively. The path to margin improvement runs through building the moat first
(\cref{ch:growth}) and then shifting the chain. \textcite{porter1985} is the primary source.

\section{Industry Structure \& Profitability}\label{sec:industry-structure}
% complexity: 3

How concentrated is the market we compete in, and what does that imply about sustainable
price levels and competitive intensity? A market with four roughly equal rivals behaves
differently from one with one dominant firm and a fringe. Concentration is not destiny ---
a concentrated market can be fiercely competitive (see: airlines) --- but it is the first
structural fact to establish. The Herfindahl--Hirschman Index translates market shares into
a single number that regulators and strategists both use.

Let firm \(i\) hold share \(s_i \in [0,1]\) (expressed as a fraction). The HHI is:
\begin{equation}\label{eq:hhi}
  \mathrm{HHI} \;=\; \sum_{i} s_i^2.
\end{equation}
Expressed in percentage-point shares (i.e.\ \(s_i\) in \([0,100]\)), the HHI ranges from
near 0 (atomistic) to \num{10000} (pure monopoly). US antitrust guidelines
\parencite{dranove7theconomics} flag markets above \num{2500} as highly concentrated and
mergers that raise HHI by more than \num{200} in a market above \num{1500}.

\cref{eq:hhi} assumes market shares proxy for competitive intensity, but two firms with
50\%\ each can collude tacitly or compete brutally. It is also sensitive to how the market
is defined: a broad definition deflates HHI and can hide local dominance. The index says
nothing about entry barriers or cost asymmetries; it is a starting point, not a conclusion.

\begin{example}
\emph{Decision: is our market concentrated enough to support price leadership, or must we
compete on value?}

Suppose the market has four players with shares
\(s = (0.30,\,0.25,\,0.20,\,0.15)\) and a competitive fringe summing to \(0.10\):
\[
  \begin{aligned}
    \mathrm{HHI} &= 0.30^2 + 0.25^2 + 0.20^2 + 0.15^2 + 0.10^2 \\
                 &= 0.090 + 0.063 + 0.040 + 0.023 + 0.010 \\
                 &= 0.226 \;\approx\; \num{2260}\text{ (pp-squared)}.
  \end{aligned}
\]
At \num{2260} the market is \emph{moderately concentrated} --- below the \num{2500}
highly-concentrated threshold. Price leadership is weak; the leading firm cannot safely
raise price without triggering share loss. A startup entering at \(s \approx 0.02\)
contributes \(0.0004\) to HHI --- negligible to the aggregate, but not to its own P\&L.
The entry is viable precisely because the concentration cliff has not been reached: no
single firm has the scale to retaliate crushingly.
\end{example}

\begin{admonition}{Warning}
Using HHI alone to conclude a market is ``competitive.'' A fragmented market of 50 firms
each with \qty{2}{\percent} share (HHI = \num{200}) can still be \emph{structurally
unattractive} if entry costs are zero and substitutes are abundant --- Five Forces says what
HHI cannot.
\end{admonition}

The Lerner index (\cref{eq:lerner}) connects HHI to price-cost margins through the
industry-average price elasticity of demand: in a symmetric Cournot oligopoly,
\(\mathcal{L} = \mathrm{HHI}/|\varepsilon|\). This links market structure directly to
pricing power and is the bridge between industrial organisation and \cref{ch:pricing}.
\parencite{mankiw2020,dranove7theconomics}.

Market concentration quantifies the arena; the value chain reveals where within it the firm
creates and captures value. Both feed the advantage analysis in
\cref{ch:advantage-and-disruption}. \textcite{mankiw2020} and \textcite{dranove7theconomics}
ground the empirical IO literature behind \cref{eq:hhi}.
